\section{The Force handler}
\label{sec:force}
Protocol~5 closes the external branch. Protocol~3 consumes the verified
closure and separately retires old programs under \cref{sec:retirement}.
The interpreter derives AUTO from accepted delivery, or applies a
previously announced MSIG Force. Both retain the entire global prefix
immediately before the triggering Espresso block's payload. Their
effective cutoff and the priority of MSIG over AUTO are unchanged.

\subsection{Export duties and block-missing windows}
\label{sec:windows}
An export duty belongs to the session responsible for the next DATA
height after $G$. Its initial authority fixes $W_{\rm start}$,
$W_{\rm run}$, $m_{\rm start}$, $q_{\min}$ and $C_{\rm auto}$.
These base-window parameters and the session quota are never rewritten.
When a duty opens in branch version $\beta$, fix its miss tolerance to
\[
 M=1+\beta
\]
and initialize $\mathit{misses}=0$. This value of $M$ remains fixed for
that duty. There is no separate retry counter. A normal handoff keeps
$\beta$; an effective Force increments it before a successor duty opens.

Let $R$ be the remaining DATA quota at a window's opening. Completion
removes the duty before $R=0$ can open another window, so every target
below is positive. At public opening timestamp $\theta_0$, create
\[
 [\theta_0,\theta_0+W_{\rm start}),\qquad
 m=\min\{m_{\rm start},R\}.
\]
This is the startup phase. After its first successful window, use
running windows of length $W_{\rm run}$ and target
$m=\min\{q_{\min},R\}$. Each window stores its phase, start, deadline,
fixed target and accepted count. The duty also stores $M$,
$\mathit{misses}$, and its total accepted DATA count $a$ for
\cref{sec:retry-slot}.

Only newly accepted DATA increments these counts. Each DATA is counted
once in its producing session. Empty DATA counts under the same quota
rule; DRAIN, heartbeats, proposal QCs, duplicate records and
future-session data behind a gap do not count. A newly opened window
starts with count zero and clips its target to the remaining quota at
that opening. Extra output is not carried as credit into a later window.

\paragraph{Settle a due window.}
At the first Espresso block whose timestamp reaches the deadline,
compare the window's count with its stored target, before reading that
block's payload.
\begin{itemize}
\item On success, set $\mathit{misses}=0$. A successful startup window
opens a running window; a successful running window opens the next
running window.
\item On failure, increment $\mathit{misses}$. If it is less than $M$,
open another base window in the \emph{same} phase. A failed startup
window therefore remains in startup. If it reaches $M$, produce one
AUTO and stop settling this duty.
\end{itemize}
Every successor base window starts at the preceding deadline, not at
the timestamp of the block that revealed the crossing. Its length is
the authority's base length for that phase. Thus successful windows
reset the consecutive-miss count, whereas a Force requires $M$
consecutive failed base windows of this duty.

If one timestamp crosses several deadlines, apply these transitions in
order until the next deadline is greater than that timestamp or an
AUTO is produced. No payload is processed between them. A terminal
DATA accepted while the version is still open cancels all remaining
checks for that session, including its miss count, even if its current
window has not expired. The half-open intervals remain significant:
DATA in a block timestamped exactly at a deadline does not repair the
window that the prelude has just settled. It may count in the next
window only if the version remains open after the prelude.

\paragraph{Batch empty-window settlement.}
Replay need not perform $M$ background callbacks. Suppose the current
window has count zero, deadline $d\le\theta$, phase length $W$, and
miss count $b<M$. With no intervening payload, all subsequently crossed
windows in this phase also have count zero and the remaining quota is
unchanged. Put
\[
 N=1+\left\lfloor\frac{\theta-d}{W}\right\rfloor,
 \qquad j=\min\{N,M-b\}.
\]
If $j=M-b$, the last required failure has deadline
$d+(j-1)W$ and causes AUTO at the current block. Otherwise set
$\mathit{misses}=b+j$ and open the window with start $d+(j-1)W$ and
deadline $d+jW$. Its count is zero. A window with a nonzero count is
settled normally first; any following zero-count run can use this
calculation. This is exactly the same transition as ordered settlement,
not a second clock or a change to a base-window length.

\paragraph{What a miss establishes.}
A failed base window establishes insufficient accepted DATA at the
specified rate. A single miss no longer necessarily closes the branch.
Receipt formation, state recovery, publication and Espresso service can
all affect the count. A missed service bound violates that liveness
premise, but the public miss does not identify the responsible party.
Neither a miss nor the resulting AUTO is automatically slashable
equivocation or proof of censorship.

\subsection{Owed service and Force selection}
\label{sec:owed}
The interpreter maintains a finite owed set $\mathcal O$ for the
current service round. Debt is discharged by accepted DATA, not by
selection, Ready, regional commitment or raw publication. When a DATA
block enters $G$, payload processing removes its producing region
from $\mathcal O$ if present. Its already assigned duty may continue
through its fixed remaining quota. An empty owed set is not refilled
while that duty continues.

At a new \emph{public} duty boundary or a pending-boundary selection,
run $\tagg{PrepareOwed}(r_*)$ for the designated target slot $r_*$:
\begin{enumerate}[label=(\roman*),itemsep=2pt]
\item Obtain the required immutable $\tagg{AUTH}_{r_*,0}$ and its
complete validation evidence. Missing evidence causes a wait at this
canonical transition; it is not evidence of zero weight or an empty set.
\item Remove members whose governance flag is OFF or whose designated
target authority has nonpositive total weight. An effective governance
OFF also removes its region when that control is applied, including an
OFF attached to a matching legacy MSIG. Stale attached configuration
has no effect.
\item If $\mathcal O$ is now empty, open a new round containing all
governance-ON regions with positive total weight under the designated
target authority. Ignore temporary AUTO cooldowns when forming this
set. If no such region exists, it remains empty and selection waits.
\end{enumerate}
Apart from this new-round initialization, an ON does not insert a
region into an unfinished round. Such a region can join the next
round. Cooldown never removes a debt, and neither CLEAR nor a new
branch version resets the owed set.

Only the canonical public interpreter calls this procedure. A
speculative REGULAR handoff may prepare a successor without changing
$\mathcal O$. Public assignment of a regular successor does call
$\tagg{PrepareOwed}$, but its selected region remains the historical
$\tagg{NextLive}$ result: REGULAR selection is not filtered by debt.

For Force, define $\tagg{NextOwed}(i,X,O)$ by searching the other
regions after source $i$ in the fixed cyclic list, restricted to
members of $O$ that are eligible in snapshot $X$ under the designated
target authority. If none exists, return $i$ itself when it is eligible
and belongs to $O$; otherwise return $\bot$. For the synthetic source
$\bot$, search the list from its first region. Governance ON,
expired cooldown and positive designated-authority weight are all
still required.

For each Espresso block, save the pre-payload snapshots
$\chi_E(q)$ and $\mathcal O_E(q)$ after its prelude. A pending forced
boundary searches from its triggering selection lower bound for the
first authenticated position at which
\[
 \tagg{NextOwed}(i_b,\chi_E(q),\mathcal O_E(q))\ne\bot,
\]
using its captured source context and designated target authority.
The first such position fixes its successor. Later payload credits or
round openings do not rewrite the saved pre-payload snapshots. A
regular witness continues to use its prescribed historical
$\tagg{NextLive}$ search, even when that region is outside the current
owed set. An ordinary OFF does not revoke an already fixed handoff.

\paragraph{Conditional AUTO cooldown.}
Before applying a new AUTO cooldown, derive its target slot by
\cref{sec:retry-slot}, run $\tagg{PrepareOwed}$, and inspect eligibility
under that target authority \emph{before} this new cooldown. Set the
source's automatic cooldown to $\theta+C_{\rm auto}$ only if another
currently eligible member of $\mathcal O$ exists. Otherwise leave its
cooldown record unchanged. This does not clear an older cooldown,
override OFF, or imply a CLEAR. It prevents this AUTO from quarantining
the last eligible source in the owed set. If existing restrictions
leave no eligible owed candidate, the boundary still waits.

\subsection{Scheduled control and two kinds of Force}
\label{sec:force-stages}
An MSIG Force binds the chain, target branch version $\beta$,
activation timestamp $\theta_a$, control sequence and authorized
attached configuration. Its L1 sequence identifies the instruction;
it supplies no retained DATA height. The committee submits it under
\cref{ass:committee,ass:monitor} and serves its finalized evidence.
Signing alone does not start the submission bound.

L1 finalization makes the instruction known. Its effect begins only at
$J_{q(\theta_a)}$. Until then, normal production, handoffs and valid
EXPORTs can continue. At activation, a target matching the entry
branch version executes; a stale target is consumed without closing
the current version or applying its attached configuration. A normal
handoff does not change $\beta$ and therefore does not cancel the
instruction. Independent STATUS records retain their own effects.

The new health policy emits MSIGs with empty attached configuration.
The interpreter nevertheless applies the authorized configuration of
a matching valid legacy record. It does not apply that configuration
early, apply a stale record's changes, or reinterpret the inherited
order of independent status changes and MSIG.

An AUTO binds the chain, old branch version, export-session pair,
fixed $M$, the sequence of consecutive misses established by replay,
the $M$th failed-window identity, and the triggering finalized header.
Its cause deadline is that last failed window's deadline; its realized
time is the triggering block's timestamp. No transaction is needed to
create the event. An optional audit record cannot execute it again.

\begin{protocol}{alg:force}{Process the prelude and close a branch}
Before the payload of authenticated finalized block $J_q$, let
$\theta=\theta_q$ and capture the entry branch version.
\begin{enumerate}[label=\textbf{\arabic*)},itemsep=3pt]
\item \textbf{Apply independent controls.} From the complete control
prefix, apply all due independent STATUS and CLEAR records not yet
applied, in $(\theta_a,c)$ order. Mark them applied. Prune debts for
governance OFF that actually takes effect. Do not add a newly ON
region to an unfinished round or refill an empty set during a
continuing duty. Later blocks with the same timestamp do not repeat
these controls.

\item \textbf{Apply a matching MSIG.} Examine due MSIGs in the same
order. Apply only the first whose target matches the entry version.
Capture the assigned duty's source context, accepted count $a$ and
index, or the stored context and target index if no duty is assigned.
Apply this record's authorized attached configuration, including any
resulting debt pruning, and close the entry branch version. Consume
every other due MSIG as stale without applying its attached changes.
If a matching MSIG closes the version, cancel its window and skip
step~3.

\item \textbf{Settle delivery windows.} If a duty exists, settle its
deadlines at or below $\theta$ under \cref{sec:windows}. Below $M$
consecutive misses, keep the branch open and retain the resulting
window state. On the $M$th miss, capture that duty's source context,
index and total accepted count $a$, and close the branch at this
position. Derive the target slot, prepare the owed set and apply the
conditional cooldown before successor selection. Missing namespace or
control evidence cannot establish a failed window. Missing target
authority evidence pauses debt, cooldown and selection computation;
it does not establish an empty candidate set or change the cutoff.

\item \textbf{Preserve the prefix and fix the retry target.} For
either closure, retain all of $G$ and its original provenance. Archive
the closed version's unconfirmed suffix and normal successors. Record
the cause, triggering header, retained anchor, captured context and
selection lower bound $q$. Under \cref{sec:retry-slot}, an assigned
duty with $a=0$ retains its current $r$ as the target slot; one with
$a>0$ targets $r+1$. Without an assigned duty, retain the stored
target slot. Increment $\beta$ exactly once. The new engine, when
entered, has fresh core state and $g=t=0$. No new-version DATA is
accepted in $J_q$ itself.

\item \textbf{Prepare, snapshot and select.} For a new boundary or
an older pending public boundary, run $\tagg{PrepareOwed}$ under its
designated target authority. Save $\chi_E(q)$ and
$\mathcal O_E(q)$ after all prelude effects. Resolve a forced boundary
by its first historical $\tagg{NextOwed}$ position; resolve a regular
boundary by its prescribed $\tagg{NextLive}$ rule, and genesis by its
synthetic selector. When a duty opens,
start its startup window at the canonical public opening timestamp,
fix $M=1+\beta$, and set $\mathit{misses}=0$ and $a=0$. If no candidate
exists, keep the anchor without a DATA window. If a duty simply
continues, still save its pre-payload snapshots but do not reset its
window, tolerance, counts or round. Then return to payload processing.
\end{enumerate}
\end{protocol}

Required selection evidence may be obtained later, but its absence
never moves a control cutoff or supplies a new opening timestamp.
Replay resumes at the same canonical transition; effects already
recorded there are not applied a second time. A Force can have a fixed
cause, new version and retained anchor while its target selection is
still unresolved.

If Espresso time jumps across several thresholds, independent
controls are applied before the matching MSIG, and MSIG takes priority
over all AUTO checks of the entry duty. Window settlement can account
for many consecutive misses, but one prelude causes at most one
effective Force. A newly assigned duty opens at the realized selection
timestamp, not at an old missed deadline, and is not settled again in
that prelude. Its positive first window prevents same-block cascading
Forces. Payload processing may then credit accepted DATA and assign
regular duties under the public boundary rules.

\paragraph{Deterministic transition.}
Given the same complete source histories and previous replay state,
the due-record order, target match, window counts, fixed $M$, miss
transitions, target slot and debt updates are fixed. Eligibility is
tested against the same designated authority, and historical
$\tagg{NextOwed}$ uses saved snapshots rather than local arrival time.
The arithmetic shortcut implements the same empty-window transitions.
Thus this prelude defines one next replay state; missing evidence
pauses its computation rather than selecting another result. This
transition statement does not supply a separate global-progress proof.

\subsection{Processing rules and control completeness}
\label{sec:control-completeness}
Control ingestion verifies a full L1 prefix and retains a pending
activation queue. A sequence number alone does not mean all ingested
records have activated. The interpreter processes each control once
in its recorded target context. A closed branch version is never
reopened by ON, repeated Force or a late regional receipt.

Every correct observer refreshes its complete finalized control view
before confirming an Espresso interval. The advance-announcement
premises ensure that all applicable controls are already available.
Missing evidence pauses replay. A new observer reconstructs the same
history; it does not reject an old instruction merely because it
learned it late.

A certifier's pure core binds its supplied L1 checkpoint, control
sequence and history commitment. Its L1 acceptance interface rejects
an attestation that omits applicable controls. A conservative
implementation requires the advertised complete control version to
match the version required by the accepting contract; a stale version
is refreshed before retry. More selective checks require an equivalent
completeness proof. One event's inclusion does not prove that no
intervening control exists.

Discovery of an omitted control that should have affected already
confirmed history violates the completeness or announcement premise.
It does not authorize rollback of $G$, a new Force at an arbitrary
later position, or a local-arrival ordering. The conditional protocol
guarantees do not specify recovery from that premise violation.
